\documentclass[12pt,a4paper]{amsart}
\usepackage[utf8]{inputenc}
\usepackage[margin=3cm]{geometry}
\usepackage[frenchb]{babel}
\geometry{letterpaper}
\usepackage[titletoc,toc,title]{appendix}
\usepackage{graphicx}
\usepackage{amssymb}
\usepackage{epstopdf}
\usepackage{mathrsfs}
\usepackage{amsmath}
\usepackage{amsthm}
\usepackage[most]{tcolorbox}
\usepackage{hyperref}
\usepackage{enumitem}
\newcommand{\ds}{\displaystyle}

\usepackage{pgfplots}
\pgfplotsset{compat=1.15}
\usepackage{mathrsfs}
\usetikzlibrary{arrows}
\pagestyle{empty}

\newcommand{\enorm}[1]{\left |\!\left |\!\left |#1\right |\!\right |\!\right |_{k}}

\newcommand{\R}{\mathbb{R}}
\newcommand{\N}{\mathbb{N}}
\newcommand{\C}{\mathbb{C}}
\newcommand{\bv}{\mathbf{v}}
\newcommand{\bx}{\mathrm{x}}
\newcommand{\cG}{c_{\rm g}}
\newcommand{\cC}{c_{\rm c}}
\newcommand{\cS}{c_{\rm s}}
\newcommand{\cN}{\LC_N}

\newcommand{\speakerone}{\bf Curely}
\newcommand{\speakertwo}{\bf Lechaux}
\newcommand{\speakerthree}{\bf Yax}
\newcommand{\speakerfour}{\bf Noel}
\newcommand{\speakerfive}{\bf Dugué}
\newcommand{\speakersix}{\bf Le Guilloux}
\newcommand{\speakerseven}{\bf Beaufort}
\newcommand{\speakereight}{\bf Morescalchi}
\newcommand{\speakernine}{\bf Chabert}
\newcommand{\speakerten}{\bf Moy}

\begin{document}
\newtheorem{theoreme}{Theorem}
\newtheorem{ex}{Example}
\newtheorem{definition}{Definition}
\newtheorem{lem}{Lemma}
\newtheorem{remarque}{Remark}
\newtheorem{exemple}{Example}
\newtheorem{proposition}{Proposition}
\newtheorem*{theo}{Theorem}
\newtheorem{corolaire}{Corollary}
\newtheorem{hyp}{Hypothesis}
\newtheorem*{rec}{Recurrence Hypothesis}

\renewcommand\Im{\mathrm{Im}~}
\renewcommand\Re{\mathrm{Re}~}
\newcommand\T{\mathbb{T}}
\newcommand\cR{\mathcal{R}}
\newcommand\cL{\mathcal{L}}
\newcommand\cX{\mathcal{X}}
\newcommand\cH{\mathcal{H}}
\newcommand\cK{\mathcal{K}}
\newcommand\Sp{\mathbb{S}}
\newcommand\Vol{\mathrm{Vol}}
\newcommand\spt{\mathrm{spt}}
\newcommand\Tr{\mathrm{Tr}}
\newcommand\WF{\mathrm{WF}}
\newcommand\Res{\mathrm{Res}}
\newcommand\n{\boldsymbol{n}}
\newcommand\dd{\mathrm{d}}
\newcommand\cD{\mathcal{D}}
\newcommand\wit{\widetilde}
\newcommand\red{\textcolor{red}}
\newcommand\hp{\mathcal{T}^{d-1}X}

\title{Rencontre des jeunes chercheureuses en analyse microlocale et géométrie}
\date{}
%\author{Organized by Yann Chaubet and Leo Tzou}
\maketitle


\section*{Schedule}
\bigskip
\vspace{0.2cm}

\begin{center}
\begin{tabular}{|c|c|c|c|}
\hline
 & Monday 02/11 & Tuesday 03/11 & Wednesday 04/11 \\
\hline

09:30--10:30 
& \speakerone
& \speakerfive
& \speakernine \\

10:30--11:00 
& Coffee
& Coffee 
& Coffee \\

11:00--12:00 
& \speakertwo
& \speakersix
&  \speakerten \\

12:00--14:00
& Lunch
& Lunch 
& Lunch  \\

14:00--15:00 
&  \speakerthree
& \speakerseven
&   \\

15:00--16:00 
& Coffee
& Coffee
&  \\

16:00--17:00 
&  \speakerfour
& \speakereight
&   \\

19:00
& 
& Dinner
&  \\

\hline
\end{tabular}
\end{center}
\vspace{2cm}

\section*{Program}

\bigskip

\subsection*{Louis Beaufort} \, New results on classical magnetic dynamics \medskip \\
\textit{Abstract: In this talk, I will introduce a new geometric framework allowing the systematic generalization of several theorems on the geodesic flow to the magnetic setting. I will discuss a few of these results, and in particular a generalization of Guillarmou-Lefeuvre's local marked length spectrum rigidity theorem, obtained in collaboration with Sebastian Munoz-Thon and Sean Richardson.}

\bigskip

\subsection*{Ambre Chabert} \, Zonal eigenmodes for the magnetic Laplacian on hyperbolic surfaces \smallskip \\
\textit{Abstract: In this talk, I will first motivate a quantization procedure for the effect of a magnetic potential on a charged particle, constrained to move on a compact Riemannian surface, by the formalism of twisted semiclassical pseudodifferential operators. When the magnetic potential is constant, and the surface is hyperbolic, I will then present how the classical dynamics reflects to the quantum level on the eigenfunctions of the quantum Hamiltonian---the magnetic Laplacian. In particular, I will introduce the critical energy level, and, for subcritical energies, construct eigenmodes with maximal sup-norm at any given point, corresponding to periodic classical orbits. This is a magnetic equivalent of the so-called zonal modes for the Laplacian on the sphere. (joint work with Thibault Lefeuvre)}

\bigskip

\subsection*{Marianne Curely} \,  Propagation of wave packets close to conical intersections \smallskip \\
\textit{Abstract: In quantum chemistry, understanding the dynamics of a system near an intersection of different energy levels is of particular interest. In particular, an important question is how transitions between these energy levels can occur. In this presentation, we study the propagation of wave packets close to conical intersections with respect to a system of two Schrödinger equations presenting a codimension 2 crossing. We focus on the dynamics that occur when the wave packets pass through a region close to the crossing, without reaching it. Our main results provide an explicit formula for the outgoing wave packet,  together with a complete description of its phase, the trajectories it follows and the drift induced by passing close to the crossing.}

\bigskip

\subsection*{Baptiste Dugué} \, Comptage orbital, mélange et chaleur en volume infini \smallskip \\
\textit{Abstract: Le problème du comptage orbital pour un groupe $\Gamma$ est l'équivalent en géométrie hyperbolique du problème du cercle de Gauss, qui consiste à compter le nombre de points à coordonnées entières dans une boule du plan euclidien et à donner un équivalent asymptotique de cette quantité. J'expliquerai comment relier cette question à celle de la vitesse de mélange du flot géodésique sur des variétés hyperboliques, et comment l'étude du noyau de la chaleur permet d'obtenir de nouveaux résultats de mélange, et donc de comptage, sur des variétés de volume infini.}

\bigskip

\subsection*{Julien Lechaux} \, Dynamics of the Magnetic Schrödinger Equation \medskip \\
\textit{Abstract: In this talk, I will discuss recent results on the long-time dynamics of the magnetic Schrödinger equation on the two-dimensional flat torus. Although the underlying classical dynamics is completely integrable, the presence of the magnetic field gives rise to regularization and equidistribution effects.}

\bigskip

\subsection*{Victor Le Guilloux} \, Average counting of geodesics on a hyperbolic surface. \medskip \\
\textit{Abstract: The moduli space is the set of hyperbolic surfaces of a given topology, up to isometries. This set is equipped with a probability measure which allows us to compute the average geometric behavior of hyperbolic surfaces in a natural way. We will first give some results of Mirzakhani on the integration of geometric random variables on the moduli space and their consequences on the average counting of simple geodesics. We will then see how to generalize Mirzakhani's integration formula to arbitrary curves and how the average counting of non-simple geodesics is related to the process of pseudo-convolution.}

\bigskip

\subsection*{Alessandro Morescalchi} \, Higher-rank Pollicott-Ruelle resonances for commuting flows \smallskip \\
\textit{Abstract: Long-time statistical properties of hyperbolic flows, namely the decay of correlations, can be analyzed in many situations with semiclassical analysis techniques, leading to a description in terms of the set of Pollicott-Ruelle resonances. In this talk we will describe how this approach can be generalized to families of commuting hyperbolic flows, arising naturally through Lie group actions, to produce a higher-rank generalization of the Pollicott-Ruelle spectrum.}

\bigskip

\subsection*{Julien Moy} \,  Spectral statistics of classically chaotic quantum systems \medskip \\
\textit{Abstract: In 1984, motivated by numerical observations on the spectrum of chaotic billiards, Bohigas, Giannoni and Schmit put forward the following conjecture: the spectrum of a quantum system whose classical analogue is chaotic should display statistical properties predicted by random matrix theory. I will first give a precise formulation of this statement, then present a heuristic argument due to Berry that supports the conjecture. Finally, I will discuss recent work that introduces randomness to the picture, in order to prove versions of the conjecture that old for "typical" quantum systems.}

\bigskip

\subsection*{Lucas Noel} \, The non-linear Schrödinger equation \smallskip \\
\textit{Abstract: The non-linear Schrödinger equation (NLS) admits special travelling waves called solitons. Due to the non-linearity, a sum of solitons is not a solution of (NLS). What we call a multi-soliton is a solution of (NLS) that behaves like a finite sum of soliton when the times grows. To prove the existence of multi-solitons, we need to understand the spectral behavior of a "linear" operator. In this talk, I would like to show how to use tools of semi-classical analysis to obtain pieces of information on the spectrum of such an operator, and what they imply.}

\bigskip

\subsection*{Arthur Yax} \, Global lower resolvent lower bound for damped waves \smallskip \\
\textit{Abstract: In this talk, we will discuss the decay rates of damped waves when the Bardos-Lebeau-Rauch-Taylor Geometric Control Condition (GCC) fails, roughly speaking "there is a geodesic that does not enter the damped region".
Since decay rates are related to resolvent growth, we  present a result of global lower resolvent estimate, as well as applications of it in various settings.}


\end{document}

